\chapter{外部资源与模型能力}

\section{外部资源}

外部资源（Resource）表示主体或调用方可以访问的知识库、文档集合、数据库、搜索服务或其他数据源。Nous Wave 保存资源身份、版本、访问范围和稳定引用，不复制成熟知识库已经提供的全文解析、同步和索引机制。

资源适配器至少提供资源描述、搜索、材料展开、版本检查和访问检查能力。搜索结果包含资源、条目、版本或修订、内容摘要或校验信息、来源区域、观察时间和访问范围。

资源搜索命中仍然属于外部材料。只有主体或认知消费者实际接触结果时才形成观察事件；只有经过目标领域形成和核验流程后才成为长期认知。

\section{稳定引用}

外部资源引用应尽可能包含提供方身份、资源身份、条目身份、版本或修订、内容摘要、来源区域坐标、获取或观察时间和访问范围。引用在资源更新后能够判断是否仍指向同一内容。

外部资源撤回、权限变化或版本变化会使依赖的派生表示和认知对象进入重新核验流程，不通过静默替换来源内容保持表面一致。

\section{模型能力}

模型和模型服务按能力合同接入。系统至少区分向量嵌入、重排、结构化提案生成、文本派生、视觉派生和语音派生等能力。

能力合同描述输入输出结构、模型或实现身份、隐私边界、可用性、成本和资源限制。所有派生内容和提案保存生产者签名，使后续重建和失效判断能够识别生成条件。

本地模型网关与远程接口在认知语义中承担相同角色。部署位置不改变权威边界；主体身份和认知状态不得绑定到某个具体模型服务。
